% コマ14: 前処理は字句解析の前 — 置換される場所・されない場所
\documentclass[border=8pt]{standalone}
% 講義資料の図(materials/sessions/figures/*.tex)共通プリアンブル
% 各図の .tex から % 講義資料の図(materials/sessions/figures/*.tex)共通プリアンブル
% 各図の .tex から % 講義資料の図(materials/sessions/figures/*.tex)共通プリアンブル
% 各図の .tex から \input{figure-preamble} で読み込む。
% 配色・フォントは latex/session-template.tex と揃える。

% 日本語(LuaLaTeX + luatexja)。等幅セル内の日本語も和文フォントで出す
\usepackage{luatexja}
\usepackage{luatexja-fontspec}
\setmainfont{Noto Sans CJK JP}
\setmainjfont{Noto Sans CJK JP}
\setmonofont{Inconsolata}[Scale=MatchLowercase]
\setmonojfont{Noto Sans CJK JP}

\usepackage{xcolor}
\definecolor{TcAccentDark}{HTML}{583B66}
\definecolor{TcAccentLight}{HTML}{EEE6F2}
\definecolor{TcAccentLighter}{HTML}{F7F3F9}
\definecolor{TcEmph}{HTML}{E64B17}
\definecolor{TcText}{HTML}{262626}
\definecolor{TcGrayText}{HTML}{737373}
\definecolor{TcGrayBg}{HTML}{F2F2F2}

\usepackage{tikz}
\usetikzlibrary{arrows.meta,positioning,calc,decorations.pathreplacing,fit,backgrounds}

\tikzset{
  % テキスト
  every node/.style={text=TcText},
  annot/.style={font=\small, text=TcText, align=left},
  gannot/.style={font=\small, text=TcGrayText, align=center},
  addr/.style={font=\ttfamily\small, text=TcGrayText},
  % メモリセル
  memcell/.style={draw=TcAccentDark, line width=0.6pt, fill=white,
                  minimum height=8.5mm, minimum width=40mm,
                  font=\ttfamily, inner sep=2pt, outer sep=0pt},
  memcell gray/.style={memcell, fill=TcGrayBg},
  memcell hi/.style={memcell, fill=TcAccentLighter},
  memcell dashed/.style={memcell, dashed, fill=none, text=TcGrayText},
  % 領域(セクション図など)
  region/.style={draw=TcAccentDark, line width=0.6pt, fill=white,
                 minimum width=48mm, align=center, inner sep=3pt, outer sep=0pt},
  % 制御フロー図のボックス
  cfgbox/.style={draw=TcAccentDark, line width=0.6pt, fill=white, rounded corners=1.5mm,
                 minimum width=28mm, minimum height=8mm, font=\small, align=center,
                 inner sep=2pt},
  % 矢印
  ptr/.style={-{Stealth[length=2.8mm]}, line width=1.1pt, color=TcEmph},
  flow/.style={-{Stealth[length=2.4mm]}, line width=0.8pt, color=TcAccentDark},
  flow back/.style={flow, dashed},
  regptr/.style={-{Stealth[length=2.8mm]}, line width=1.1pt, color=TcAccentDark},
  % 補助
  bracestyle/.style={decorate, decoration={brace, amplitude=4pt}, line width=0.8pt,
                     color=TcAccentDark},
}
 で読み込む。
% 配色・フォントは latex/session-template.tex と揃える。

% 日本語(LuaLaTeX + luatexja)。等幅セル内の日本語も和文フォントで出す
\usepackage{luatexja}
\usepackage{luatexja-fontspec}
\setmainfont{Noto Sans CJK JP}
\setmainjfont{Noto Sans CJK JP}
\setmonofont{Inconsolata}[Scale=MatchLowercase]
\setmonojfont{Noto Sans CJK JP}

\usepackage{xcolor}
\definecolor{TcAccentDark}{HTML}{583B66}
\definecolor{TcAccentLight}{HTML}{EEE6F2}
\definecolor{TcAccentLighter}{HTML}{F7F3F9}
\definecolor{TcEmph}{HTML}{E64B17}
\definecolor{TcText}{HTML}{262626}
\definecolor{TcGrayText}{HTML}{737373}
\definecolor{TcGrayBg}{HTML}{F2F2F2}

\usepackage{tikz}
\usetikzlibrary{arrows.meta,positioning,calc,decorations.pathreplacing,fit,backgrounds}

\tikzset{
  % テキスト
  every node/.style={text=TcText},
  annot/.style={font=\small, text=TcText, align=left},
  gannot/.style={font=\small, text=TcGrayText, align=center},
  addr/.style={font=\ttfamily\small, text=TcGrayText},
  % メモリセル
  memcell/.style={draw=TcAccentDark, line width=0.6pt, fill=white,
                  minimum height=8.5mm, minimum width=40mm,
                  font=\ttfamily, inner sep=2pt, outer sep=0pt},
  memcell gray/.style={memcell, fill=TcGrayBg},
  memcell hi/.style={memcell, fill=TcAccentLighter},
  memcell dashed/.style={memcell, dashed, fill=none, text=TcGrayText},
  % 領域(セクション図など)
  region/.style={draw=TcAccentDark, line width=0.6pt, fill=white,
                 minimum width=48mm, align=center, inner sep=3pt, outer sep=0pt},
  % 制御フロー図のボックス
  cfgbox/.style={draw=TcAccentDark, line width=0.6pt, fill=white, rounded corners=1.5mm,
                 minimum width=28mm, minimum height=8mm, font=\small, align=center,
                 inner sep=2pt},
  % 矢印
  ptr/.style={-{Stealth[length=2.8mm]}, line width=1.1pt, color=TcEmph},
  flow/.style={-{Stealth[length=2.4mm]}, line width=0.8pt, color=TcAccentDark},
  flow back/.style={flow, dashed},
  regptr/.style={-{Stealth[length=2.8mm]}, line width=1.1pt, color=TcAccentDark},
  % 補助
  bracestyle/.style={decorate, decoration={brace, amplitude=4pt}, line width=0.8pt,
                     color=TcAccentDark},
}
 で読み込む。
% 配色・フォントは latex/session-template.tex と揃える。

% 日本語(LuaLaTeX + luatexja)。等幅セル内の日本語も和文フォントで出す
\usepackage{luatexja}
\usepackage{luatexja-fontspec}
\setmainfont{Noto Sans CJK JP}
\setmainjfont{Noto Sans CJK JP}
\setmonofont{Inconsolata}[Scale=MatchLowercase]
\setmonojfont{Noto Sans CJK JP}

\usepackage{xcolor}
\definecolor{TcAccentDark}{HTML}{583B66}
\definecolor{TcAccentLight}{HTML}{EEE6F2}
\definecolor{TcAccentLighter}{HTML}{F7F3F9}
\definecolor{TcEmph}{HTML}{E64B17}
\definecolor{TcText}{HTML}{262626}
\definecolor{TcGrayText}{HTML}{737373}
\definecolor{TcGrayBg}{HTML}{F2F2F2}

\usepackage{tikz}
\usetikzlibrary{arrows.meta,positioning,calc,decorations.pathreplacing,fit,backgrounds}

\tikzset{
  % テキスト
  every node/.style={text=TcText},
  annot/.style={font=\small, text=TcText, align=left},
  gannot/.style={font=\small, text=TcGrayText, align=center},
  addr/.style={font=\ttfamily\small, text=TcGrayText},
  % メモリセル
  memcell/.style={draw=TcAccentDark, line width=0.6pt, fill=white,
                  minimum height=8.5mm, minimum width=40mm,
                  font=\ttfamily, inner sep=2pt, outer sep=0pt},
  memcell gray/.style={memcell, fill=TcGrayBg},
  memcell hi/.style={memcell, fill=TcAccentLighter},
  memcell dashed/.style={memcell, dashed, fill=none, text=TcGrayText},
  % 領域(セクション図など)
  region/.style={draw=TcAccentDark, line width=0.6pt, fill=white,
                 minimum width=48mm, align=center, inner sep=3pt, outer sep=0pt},
  % 制御フロー図のボックス
  cfgbox/.style={draw=TcAccentDark, line width=0.6pt, fill=white, rounded corners=1.5mm,
                 minimum width=28mm, minimum height=8mm, font=\small, align=center,
                 inner sep=2pt},
  % 矢印
  ptr/.style={-{Stealth[length=2.8mm]}, line width=1.1pt, color=TcEmph},
  flow/.style={-{Stealth[length=2.4mm]}, line width=0.8pt, color=TcAccentDark},
  flow back/.style={flow, dashed},
  regptr/.style={-{Stealth[length=2.8mm]}, line width=1.1pt, color=TcAccentDark},
  % 補助
  bracestyle/.style={decorate, decoration={brace, amplitude=4pt}, line width=0.8pt,
                     color=TcAccentDark},
}

\begin{document}
\begin{tikzpicture}[node distance=0mm,
    stg/.style={cfgbox, minimum width=22mm, minimum height=8mm, font=\small},
    ftxt/.style={font=\ttfamily\small, inner sep=1.5pt},
    srcline/.style={font=\ttfamily\footnotesize, text=TcText, anchor=west},
    okmark/.style={annot, text=TcEmph, anchor=west},
    note/.style={annot, text=TcGrayText, anchor=west}]

  % ================ 上段: パイプライン上の位置 ================
  \node[gannot, anchor=west] at (-74mm, 14mm) {前処理はパイプラインのどこで効くか};

  \node[ftxt, anchor=west] (f1) at (-74mm, 0mm) {ソース文字列};
  \node[stg, right=4mm of f1, fill=TcAccentLighter, line width=1.1pt] (s1) {preprocess};
  \node[ftxt, right=4mm of s1] (f2) {展開済み文字列};
  \node[stg, right=4mm of f2] (s2) {Lexer};
  \node[ftxt, right=4mm of s2] (f3) {トークン列};
  \node[stg, right=4mm of f3] (s3) {Parser};
  \node[ftxt, right=4mm of s3] (f4) {AST};

  \foreach \a/\b in {f1/s1, s1/f2, f2/s2, s2/f3, f3/s3, s3/f4} {
    \draw[flow] (\a.east) -- (\b.west);
  }

  \node[annot, text=TcEmph, anchor=north, align=center] at ($(s1.south)+(0mm,-2mm)$)
      {ここで {\ttfamily \#include} と {\ttfamily \#define} が消える};
  \node[gannot, anchor=north, align=center] at ($(s2.south)+(0mm,-2mm)$)
      {Lexer は指令行を\\一度も見ない};

  % ================ 下段: 置換される場所・されない場所 ================
  \node[gannot, anchor=west] at (-74mm, -22mm) {同じ綴りの {\ttfamily BASE} でも、置換されるのは識別子トークンだけ};

  \node[region, minimum width=54mm, minimum height=34mm, anchor=north west] (src) at (-74mm, -32mm) {};
  \node[gannot, anchor=south west] at ($(src.north west)+(1mm,0.5mm)$) {前処理前のソース};
  \node[srcline] at ($(src.north west)+(3mm,-6mm)$)  {\#define BASE 2};
  \node[srcline] at ($(src.north west)+(3mm,-14mm)$) {return BASE + 1;};
  \node[srcline] at ($(src.north west)+(3mm,-22mm)$) {printf("BASE");};
  \node[srcline] at ($(src.north west)+(3mm,-30mm)$) {// BASE};

  \node[region, minimum width=54mm, minimum height=34mm, anchor=north west] (dst) at (4mm, -32mm) {};
  \node[gannot, anchor=south west] at ($(dst.north west)+(1mm,0.5mm)$) {前処理後（Lexer が受け取る文字列）};
  \node[srcline, text=TcGrayText] at ($(dst.north west)+(3mm,-6mm)$)  {（行ごと消える）};
  \node[srcline] at ($(dst.north west)+(3mm,-14mm)$) {return {\color{TcEmph}2} + 1;};
  \node[srcline] at ($(dst.north west)+(3mm,-22mm)$) {printf("BASE");};
  \node[srcline] at ($(dst.north west)+(3mm,-30mm)$) {// BASE};

  \draw[ptr] (src.east) -- node[gannot, above=0.5mm] {preprocess} (dst.west);

  \node[note] at (62mm, -38mm) {指令行そのものが取り除かれる};
  \node[okmark] at (62mm, -46mm) {○ 識別子トークンなので {\ttfamily 2} に置換される};
  \node[note] at (62mm, -54mm) {× 文字列リテラルの中身はトークンではない};
  \node[note] at (62mm, -62mm) {× コメントの中身はトークンではない};

  % ---- 補足 ----
  \node[gannot, anchor=north, align=center] at (10mm, -74mm)
      {置換は文字列のうえで起き、その結果が Lexer に渡る。だから AST に {\ttfamily BASE} は残らない。\\
       逆に {\ttfamily "BASE"} は文字列リテラルの中身のまま {\ttfamily .data} に出る};
\end{tikzpicture}
\end{document}
